% ============================================================
%  The cast -- twelve dossiers
% ============================================================
\rsection{The cast}

\begin{roadmap}[Twelve functions carry the whole document]
Each dossier gives the definition, the property the function has, the
property it lacks, and the arrow it kills. Several are Rudin's own and are
annotated in the main edition; those references are given.

\medskip
\raggedright
\begin{tabular}{@{}l l l@{}}
\toprule
& \textbf{function} & \textbf{kills the claim that\ldots} \\
\midrule
E1  & removable spike            & limit $\Rightarrow$ continuous \\
E2  & $|x|$                      & continuous $\Rightarrow$ differentiable \\
E3  & $x^2\sin(1/x)$             & differentiable $\Rightarrow$ $\mathscr{C}^1$ \\
E4  & $\sin(1/x)$                & bounded $\Rightarrow$ has a limit \\
E5  & step function              & integrable $\Rightarrow$ continuous \\
E6  & Thomae                     & integrable $\Rightarrow$ disc.\ set nowhere dense \\
E7  & Dirichlet                  & bounded $\Rightarrow$ integrable \\
E8  & Weierstrass                & continuous $\Rightarrow$ differentiable anywhere \\
E9  & $x^n$ on $[0,1]$           & pointwise preserves continuity \\
E10 & $\sin(nx)/\sqrt n$         & uniform preserves differentiability \\
E11 & tall spikes                & pointwise preserves $\int$ \\
E12 & Volterra                   & derivative $\Rightarrow$ integrable \\
\bottomrule
\end{tabular}
\end{roadmap}

\ritem{E1}{Removable spike}
$f(x)=0$ for $x \ne 0$, $f(0)=1$. Has limit $0$ at $0$; not continuous
there, since the limit is not the value. The cheapest possible separation,
and the reason Rudin's limit is \emph{deleted} (S4.2).

\ritem{E2}{Absolute value}
$f(x)=|x|$. Continuous everywhere; at $0$ the difference quotient is
$|h|/h = \pm1$ according to the side, so the one-sided derivatives are $-1$
and $+1$ and $f'(0)$ does not exist. Note that $f$ is Lipschitz, so even
that stronger condition does not give differentiability.

\ritem{E3}{$x^2\sin(1/x)$}
$f(x) = x^2\sin(1/x)$ for $x \ne 0$, $f(0)=0$. Then
$|f(x)| \le x^2$ forces $f'(0) = 0$, and for $x \ne 0$
\[ f'(x) = 2x\sin\frac1x - \cos\frac1x , \]
which has no limit as $x \to 0$ --- the second term oscillates between $-1$
and $1$. So $f$ is differentiable everywhere with $f'$ discontinuous at $0$;
the discontinuity is of the second kind, as T4 requires. This function is
also the building block of E12.

\ritem{E4}{$\sin(1/x)$}
Bounded by $1$; no limit at $0$, since $x_n = 1/(2n\pi)$ and
$y_n=1/(2n\pi + \pi/2)$ both tend to $0$ while $f(x_n)=0$ and $f(y_n)=1$.
Rudin's 4.27(d); the oscillatory case of the four-way taxonomy at S4.19, and
the reason $\omega_f(0)=2$.

\ritem{E5}{Step function}
$f = 0$ on $[0,1)$, $f=1$ on $[1,2]$. Integrable --- one discontinuity, a set
of measure zero (T10) --- and not continuous. By T4 it is also \emph{not} a
derivative, since it has a jump. So E5 separates ``integrable'' from ``has an
antiderivative'' in one direction, and E12 separates them in the other.

\ritem{E6}{Thomae}
$f(p/q) = 1/q$ in lowest terms, $f = 0$ at irrationals. Continuous at every
irrational and discontinuous at every rational. Integrable by T10, because
$\Q$ is countable and therefore of measure zero --- with $\int_0^1 f = 0$.
Rudin's Exercise 4.18; S4.5 in the Chapter 4 appendix proves the companion
fact that no function is continuous exactly on $\Q$, so Thomae has no mirror
image.

\ritem{E7}{Dirichlet}
$\mathbf 1_{\Q}$. Bounded, discontinuous at every point, so
$\omega_f \equiv 1$ and by T10 it is not integrable --- directly, every upper
sum is $1$ and every lower sum is $0$. It is the pointwise limit of the
integrable indicators of finite rational sets (T18) but, by Baire, \emph{not}
of any sequence of continuous functions (S2.13).

\ritem{E8}{Weierstrass}
$f(x) = \sum_{k=0}^{\infty} a^k\cos(b^k\pi x)$ with $0<a<1$, $b$ an odd
integer, $ab > 1 + \tfrac32\pi$. The series converges uniformly by the
$M$-test, so $f$ is continuous (T14); it is differentiable at no point.
Rudin's 7.18. Baire's theorem (S2.13) upgrades the example to a statement
about the typical case: the continuous functions differentiable somewhere
form a set of first category in $\mathscr{C}[a,b]$, so nowhere
differentiability is the rule and E2 the exception.

\ritem{E9}{$x^n$ on $[0,1]$}
Each $f_n$ is continuous; the pointwise limit is $0$ on $[0,1)$ and $1$ at
$1$. Convergence is not uniform, since
$\sup_{[0,1]}|f_n - f| = 1$ for every $n$ --- take $x$ just below $1$.
Kills ``pointwise preserves continuity'' (T18), and is the standard picture
of a limit escaping an $\varepsilon$-tube.

\ritem{E10}{$\sin(nx)/\sqrt n$}
$\sup|f_n| = 1/\sqrt n \to 0$, so $f_n \to 0$ uniformly; but
$f_n'(x) = \sqrt n\cos(nx)$ diverges everywhere. Kills ``uniform preserves
differentiability'' (T16). The mechanism is worth naming: differentiating
multiplies the amplitude by the frequency, and uniform convergence bounds
only the amplitude.

\ritem{E11}{Tall spikes}
$f_n = n \cdot \mathbf 1_{(0,1/n)}$. Each is integrable with
$\int_0^1 f_n = 1$; $f_n \to 0$ pointwise. Kills
``pointwise gives $\int \lim = \lim \int$'' (T18). The functions are not
uniformly bounded, which is exactly the hypothesis the dominated convergence
theorem (11.32) would need.

\ritem{E12}{Volterra}
Built on a fat Cantor set $K$ --- nowhere dense, but not of measure zero ---
by planting a tapered copy of E3 in each removed interval and setting the
function to $0$ on $K$. The result is differentiable everywhere with bounded
derivative, and the derivative oscillates at every point of $K$, so
$\Omega_1 \supset K$ and by T10 the derivative is not integrable (T12).
The single most informative example in the document: it shows the two halves
of the Fundamental Theorem are independent, and it is why 6.21 must
hypothesise integrability rather than deduce it.

\begin{recap}[What to carry away]
\raggedright
\textbf{One.} The local chain is strict at every link, and the witnesses are
cheap --- E1, E2, E3. Nothing subtle is needed to separate limit,
continuity, differentiability, and continuous differentiability.

\smallskip
\textbf{Two.} Integrability is not part of that chain. It is global, it comes
from compactness (T8), and it has an exact characterisation in terms of the
size of the discontinuity set (T10) --- for which the oscillation machinery
of S2.5 and S4.5 is precisely the right language.

\smallskip
\textbf{Three.} The interesting failures are all failures of \emph{exchange}:
$\lim$ against $\int$, $\lim$ against $\lim$, $\lim$ against $\frac{d}{dx}$.
Uniform convergence fixes the first two and cannot fix the third, because
differentiation amplifies exactly what uniformity controls.

\smallskip
\textbf{Four.} Where an implication fails, ask what hypothesis would restore
it. Pointwise $\to$ uniform restores continuity and the integral; uniform on
the functions $\to$ uniform on the derivatives restores differentiation
(T17); Riemann $\to$ Lebesgue restores the integral exchange under a
domination hypothesis instead of a uniformity one. Every gap in the lattice
is an invitation of this kind.
\end{recap}
